% ----------------------------------------------------------------------
% common text abbreviations
\newcommand{\eg}{\mbox{\it e.g.}}
\newcommand{\ie}{\mbox{\it i.e.}}
\newcommand{\etal}{\mbox{\it et al.}}
\newcommand{\apriori}{{\it a priori}}
\newcommand{\aposteriori}{{\it a posteriori}}
\newcommand{\wloss}{without loss of generality}
\newcommand{\Wloss}{Without loss of generality}

% ----------------------------------------------------------------------
% common symbols

% implication
\newcommand{\imp}{\Rightarrow}                  % implication (math mode)
\newcommand{\iimp}{\Longrightarrow}             % long implication (math mode)
\renewcommand{\implies}{$\imp\ $}		% '=>' (outside math mode)
\newcommand{\pmi}{\Leftarrow}			% <= (opposite of \imp) 

% spacing
\newcommand{\sep}{\hspace{2em}} 
\newcommand{\ssep}{\hspace{1em}}
\newcommand{\void}[1]{}				% comment out the argument
\newcommand{\todo}[1]{{\color{blue}[#1]\marginpar{\mbox{\Huge !}}}}

%Here are some user-defined notations

\newcommand{\RR}{\mathbf R}
\newcommand{\CC}{\mathbf C}
\newcommand{\ZZ}{\mathbf Z}
\newcommand{\ZZn}[1]{\ZZ/{#1}\ZZ}
\newcommand{\QQ}{\mathbf Q}
\newcommand{\rr}{\mathbb R}
\newcommand{\cc}{\mathbb C}
\newcommand{\zz}{\mathbb Z}
\newcommand{\zzn}[1]{\zz/{#1}\zz}
\newcommand{\qq}{\mathbb Q}
\newcommand{\nn}{\mathbb N}
\newcommand{\calM}{\mathcal M}
\newcommand{\dd}{{\rm d}}
\newcommand{\sm}{\setminus} 
\newcommand{\set}[1]{\mathcal{#1}}
\newcommand{\Lap}[1]{\mathcal{L}}

\newcommand{\thisway}{$(\Longrightarrow)$}
\newcommand{\thatway}{$(\Longleftarrow)$}


\newcommand{\Prob}{\operatorname {Prob}}


% ----------------------------------------------------------------------
% theorem definitions
% numbered by default: 
% 	theorem, lemma, proposition, corollary, example, principle
% unnumbered by default: 
%	definition, openproblem, remark, remarks, convention
%
% get unnumbered versions by un-lemma etc., numbered by n-definition etc. 

% \theoremstyle{plain}
% \newtheorem{n-lemma}{Lemma}[chapter]
% \newtheorem{n-theorem}[n-lemma]{Theorem}
% \newtheorem{n-proposition}[n-lemma]{Proposition}
% \newtheorem{n-corollary}[n-lemma]{Corollary} 
% \newtheorem{n-principle}[n-lemma]{Principle}
% \newtheorem{n-claim}[n-lemma]{Claim}

% \newtheorem*{un-lemma}{Lemma}
% \newtheorem*{un-theorem}{Theorem}
% \newtheorem*{un-proposition}{Proposition}
% \newtheorem*{un-corollary}{Corollary} 
% \newtheorem*{un-principle}{Principle}
% \newtheorem*{un-claim}{Claim}

% \theoremstyle{definition}
% \newtheorem{n-definition}[n-lemma]{Definition}
% \newtheorem{n-openproblem}[n-lemma]{Open Problem}
% \newtheorem{bn-remark}[n-lemma]{Remark}		% bold remark
% \newtheorem*{un-definition}{Definition}
% \newtheorem*{un-openproblem}{Open Problem}
% \newtheorem*{bun-remark}{Remark}		% bold unnumbered remark
% \newtheorem*{un-problem}{Problem}
% \newtheorem{n-problem}[n-lemma]{Problem}


% \theoremstyle{remark}
% \newtheorem{n-example}[n-lemma]{Example}
% \newtheorem{n-remark}[n-lemma]{Remark}
% \newtheorem{n-remarks}[n-lemma]{Remarks}
% \newtheorem{n-convention}[n-lemma]{Convention}
% \newtheorem*{un-example}{Example}
% \newtheorem*{un-remark}{Remark}
% \newtheorem*{un-remarks}{Remarks}
% \newtheorem*{un-convention}{Convention}

% % defaults numbered/unnumbered
% \newenvironment{theorem}{\begin{n-theorem}}{\end{n-theorem}}
% \newenvironment{lemma}{\begin{n-lemma}}{\end{n-lemma}}
% \newenvironment{proposition}{\begin{n-proposition}}{\end{n-proposition}}
% \newenvironment{corollary}{\begin{n-corollary}}{\end{n-corollary}}
% \newenvironment{definition}{\begin{n-definition}}{\end{n-definition}}
% \newenvironment{example}{\begin{n-example}}{\end{n-example}}
% \newenvironment{openproblem}{\begin{un-openproblem}}{\end{un-openproblem}}
% \newenvironment{principle}{\begin{n-principle}}{\end{n-principle}}
% \newenvironment{remark}{\begin{n-remark}}{\end{n-remark}}
% \newenvironment{remarks}{\begin{un-remarks}}{\end{un-remarks}}
% \newenvironment{convention}{\begin{un-convention}}{\end{un-convention}}
% \newenvironment{claim}{\begin{un-claim}}{\end{un-claim}}
% \newenvironment{problem}{\begin{n-problem}}
% {\end{n-problem}}
% % proofs
\renewenvironment{proof}
{\noindent{\it Proof:}\hspace{.5em}\rm}
{\hfill$\square$

\vspace{0.1em}}

\newtheorem{assumption}{Assumption}

\newenvironment{proofof}[1]			% loose proofs
	{\begin{trivlist}\item[]{\it Proof of #1:\hspace{.5em}}\rm}
	{\hfill$\square$
	\end{trivlist}}
	%{\hfill\qedsymbol
	%\end{trivlist}}
\newenvironment{proofsketch}			% proof sketch
	{\begin{trivlist}\item[]{\it Proof sketch:\hspace{.5em}}\rm}
	{\hfill$\square$
	\end{trivlist}}

% \renewcommand{\qed}{} 		% turn off automatic AMS-style \qed'ing.
% \newcommand{\eotsymbol}{{\large $\Box$}} 	% can be redefined
% \newcommand{\eot}{{\hspace*{\fill}\eotsymbol}}
% \newcommand{\eottwo}{\vspace{-5ex}\par\eot}	% \eot after math display

% ----------------------------------------------------------------------
% case distinctions: label is flush left. 
% caselist causes no indentation. May use \case for automatic numbering
% itemlist causes indentation. 
% nienumerate is like enumerate but flush left

% \newcounter{casecounter}
% \newlength{\normalparindent}
% \setlength{\normalparindent}{\parindent}
% \newenvironment{caselist}
% 	{\begin{list}{{\bf Case \arabic{casecounter}:}}
% 	 {\usecounter{casecounter}
% 		\setlength{\leftmargin}{0ex}
% 		\setlength{\labelwidth}{0ex} 		 
% 		\setlength{\labelsep}{2ex}
% 		\setlength{\itemindent}{\labelsep}	
% 		\setlength{\parsep}{0ex}
% 		\setlength{\listparindent}{\normalparindent}
% 	}}	
% 	{\end{list}}
% \newcommand{\case}{\item}
% \newcommand{\caseimp}{\case[$\imp$:]}
% \newcommand{\casepmi}{\case[$\pmi$:]}
% \newcommand{\impl}[2]{$\text{#1}\imp\text{#2}$}
% \newcommand{\caseimpl}[2]{\case[\impl{#1}{#2}:]}
% \newenvironment{itemlist}{\begin{list}{\undefinedcase}
%  	{	\setlength{\leftmargin}{2.5ex}
% 		\setlength{\labelwidth}{1.5ex}
%  		\setlength{\labelsep}{1ex}
% 		\setlength{\itemindent}{0ex}}}
%  	{\end{list}}

% \newcounter{nicounter}
% \newenvironment{nienumerate} % enumerate but no indentation
% 	{\begin{list}{\arabic{nicounter}.}
% {\usecounter{nicounter} 		 \setlength{\leftmargin}{0ex}
% \setlength{\labelwidth}{0ex} 		 \setlength{\labelsep}{1ex}
% \setlength{\itemindent}{\labelsep} 	}	} 	{\end{list}}


%This numbers equations by section
\numberwithin{equation}{section}


\newcommand{\assign}{\leftarrow}
\newcommand{\True}{\operatorname{\mathbf{True}}}
\newcommand{\False}{\operatorname{\mathbf{False}}}

\newcommand{\U}{{\mathcal{U}}}
\renewcommand{\L}{{\mathcal{L}}}
\newcommand{\upperbd}[2]{\beta^+(\U_{#1, #2})}
\newcommand{\lowerbd}[2]{\beta^-(\L_{#1, #2})}
% \newcommand{\upperbdwalk}[2]{\mathfrak{UW}_{#1,#2}}
% \newcommand{\lowerbdwalk}[2]{\mathfrak{LW}_{#1,#2}}
% \newcommand{\upperbdwalks}{\upperbdwalk}
% \newcommand{\lowerbdwalks}{\lowerbdwalk}
\newcommand{\upperbdpath}[2]{\mathcal{U}_{#1,#2}}
\newcommand{\lowerbdpath}[2]{\mathcal{L}_{#1,#2}}
\newcommand{\upperbdpaths}{\upperbdpath}
\newcommand{\lowerbdpaths}{\lowerbdpath}
% \newcommand{\mupperbdpath}[2]{\mathfrak{MUP}_{#1,#2}}
% \newcommand{\mlowerbdpath}[2]{\mathfrak{MLP}_{#1,#2}}
% \newcommand{\mupperbdpaths}{\mupperbdpath}
% \newcommand{\mlowerbdpaths}{\mlowerbdpath}

\newcommand{\robinson}[1]{\mathscr{#1}}
\newcommand{\bm}{\boldsymbol}
\newcommand{\minimal}{\operatorname{minimal}}
\newcommand{\maximal}{\operatorname{maximal}}
\newcommand{\dist}[2]{\Pi(#2)-\Pi(#1)}%{\operatorname{dist}}
\newcommand{\nedge}[0]{null-edge}
\newcommand{\embedding}{uniform embedding}

\newcommand{\away}[0]{away-edge}
\newcommand{\within}[0]{within-edge}

% \renewcommand{\leq}{\leqslant}		
% \renewcommand{\geq}{\geqslant}	
% \renewcommand{\preceq}{\preccurlyeq}
% \renewcommand{\succeq}{\succcurlyeq}

\newcommand{\diagincmat}{Robinson matrix}

\newcommand{\path}[1]{\langle#1\rangle}	
\newcommand{\bma}{\bm a}
\newcommand{\bmb}{\bm b}
\newcommand{\bmc}{\bm c}
\newcommand{\bmd}{\bm d}
\newcommand{\bmx}{\bm x}
\newcommand{\bmy}{\bm y}
\newcommand{\bmw}{\bm w}


\newcommand{\calP}{\mathcal P}
\newcommand{\calS}{\mathcal S}
\newcommand{\calH}{\mathcal H}

\newcommand{\st}{\operatorname{s.t.\ }}

\newcommand{\upperbounds}{\operatorname{UpperBounds}}
\newcommand{\lowerbounds}{\operatorname{LowerBounds}}

\newcommand{\edge}[2]{\{#1, #2\}}
\newcommand{\reverse}[1]{#1^{\leftarrow}}
\newcommand{\rev}{^{\leftarrow}}
\newcommand{\bmD}{\mathbb D^k}